% ==============================================================================
% CHƯƠNG 2: TỔNG QUAN TÀI LIỆU
% ==============================================================================

\chapter{Tổng quan tài liệu}\label{ch:tong_quan}

\startcontents[chapters]
\printmyminitoc{
Chương này trình bày cơ sở lý thuyết nền tảng, mô hình toán học liên quan và tổng hợp các công trình nghiên cứu trước đây \cite{sample_book2022}.
}

\section{Cơ sở lý thuyết và Công thức toán học}\label{sec:ch02_co_so}

Các công thức toán học trong bài viết được đánh số tự động theo chương bằng môi trường \texttt{equation}. Ví dụ hàm mục tiêu chuẩn hóa:
\begin{equation}\label{eq:sample_loss}
  \mathcal{L}(\theta) = \frac{1}{N} \sum_{i=1}^{N} \left( y_i - f(x_i; \theta) \right)^2
\end{equation}
trong đó $N$ là số lượng mẫu dữ liệu, $y_i$ là giá trị thực tế và $f(x_i; \theta)$ là giá trị dự báo từ mô hình tại tham số $\theta$.

\section{Tổng hợp các công trình liên quan}\label{sec:ch02_lien_quan}

So sánh, phân tích ưu điểm và hạn chế của các nghiên cứu trước đây \cite{sample_article2024, sample_book2022} nhằm chỉ ra khoảng trống nghiên cứu mà đề tài sẽ tập trung giải quyết.
